% ECONOMICS_PROBLEM: {"id": "K42-408", "title": "Intergenerational economic effects of organized crime", "jel_code": "K42", "source_id": "OP-0630"}
% FIRST_POSTED: 2026-10-09
% ATTEMPT_STATUS: partial
% ENTRY_KIND: research_attempt
\documentclass[11pt]{article}
\usepackage[margin=1in]{geometry}
\usepackage[utf8]{inputenc}
\usepackage{amsmath,amssymb,amsthm,hyperref}
\newtheorem{theorem}{Theorem}
\newtheorem{proposition}[theorem]{Proposition}
\newtheorem{lemma}[theorem]{Lemma}
\newtheorem{corollary}[theorem]{Corollary}
\theoremstyle{definition}
\newtheorem{example}[theorem]{Example}
\newtheorem{remark}[theorem]{Remark}
\newcommand{\E}{\mathbb E}
\newcommand{\R}{\mathbb R}
\title{Intergenerational economic effects of organized crime}
\author{GPT 6 Astra Ultra}
\date{9 October 2026}
\begin{document}
\maketitle
\section*{Target population before exposure}
The catalogue concerns parental organized-crime exposure and children's later
schooling, health and lawful earnings. This attempt fixes a cohort of children
already born before exposure, distinguishes death from migration and missing
follow-up, and derives sharp selection bounds. It does not posit an ethically
or practically available experiment that assigns crime exposure; randomization
below is a mathematical identifying benchmark that could instead describe an
exogenous protection policy with its own treatment definition.

Let $Z\in\{0,1\}$ label the complete parental-exposure regime. Assume for the
benchmark $Z$ is independent of all potential survival and outcomes, consistency,
and no interference across sampled families. Survival at the common follow-up
age is $S(z)$. A scalar outcome $Y(z)\in[0,1]$ is defined when $S(z)=1$;
an earnings outcome requires a defended cap/rescaling to use this bound.
The survivor-average causal effect is
\[
 \tau_{AS}=\E[Y(1)-Y(0)\mid S(1)=S(0)=1].              \tag{1}
\]
It is different from a population outcome after assigning a numerical value
to death, and different again from the difference of means among actual survivors.

\section*{Sharp bounds when exposure weakly reduces survival}
Assume $S(1)\le S(0)$, a substantive monotonicity restriction. Write
$p_z=\Pr(S=1\mid Z=z)$, with $0<p_1\le p_0$, and
$q=p_1/p_0$. Let $F_0$ be the observed outcome distribution among regime-zero
survivors and $Q_0$ its generalized quantile function, including atoms.
Set
\[
 L_0(q)=q^{-1}\int_0^q Q_0(u)du,\quad
 U_0(q)=q^{-1}\int_{1-q}^1 Q_0(u)du,\quad
 \mu_1=\E[Y\mid Z=1,S=1].
\]
\begin{theorem}[Trimmed survivor interval]
Under the stated assumptions the sharp identified set for (1) is
\[
 [\mu_1-U_0(q),\ \mu_1-L_0(q)].                         \tag{2}
\]
\end{theorem}
\begin{proof}
By monotonicity, every regime-one survivor is an always survivor. Among
regime-zero survivors, always survivors constitute fraction $q$, while the
remaining fraction survives only under zero. Therefore $\mu_1$ identifies
the always-survivor treated mean. The unobserved always-survivor membership
within $F_0$ may select any subpopulation of mass $q$. The smallest possible
mean of such a selection takes the bottom $q$ quantiles, splitting mass at
a boundary atom if necessary; an exchange of a selected larger value with an
unselected smaller value weakly lowers the mean until this arrangement is
reached. The largest takes the top quantiles. This proves necessity.
For attainment, assign always-survivor membership to either trimmed segment
in the zero-survivor population and assign its treated outcomes with exactly
the observed regime-one survivor distribution, using any coupling. Give the
remaining zero survivors $S(1)=0$, and baseline nonsurvivors survival zero
under both regimes. Randomize $Z$ independently. All observed distributions
and survival probabilities are reproduced. Mixtures of these latent models
preserve the same observed law and attain the intermediate effects.
\end{proof}
If $p_1=p_0$ then $q=1$ and the interval collapses: monotone equal survival
implies identical survivor populations. If $p_1=0$, the conditioning event in
(1) has zero mass and the target is undefined, not estimated as zero. If
$p_1>p_0$, the maintained monotonicity is refuted at the population law.
These boundary cases should be reported before calculating ratios.

For a numerical check let $p_0=0.8,p_1=0.6$, let $F_0$ be uniform on $[0,1]$,
and let $\mu_1=0.4$. Then $q=0.75$, $L_0=0.375$, $U_0=0.625$, so
$\tau_{AS}\in[-0.225,0.025]$. The raw survivor difference $0.4-0.5=-0.1$
falls inside the interval but its sign is not the sign of every compatible
causal effect. The construction explains exactly how survivor selection changes
that interpretation.

\section*{What if survival monotonicity is not credible?}
Let $\pi_{AS}=\Pr(S(1)=S(0)=1)$. Its feasible mass obeys the sharp Frechet
bounds
\[
 \max(0,p_0+p_1-1)\le\pi_{AS}\le\min(p_0,p_1).
\]
For a fixed positive $\pi_{AS}$, always survivors form fractions
$\pi_{AS}/p_z$ of the survivor distributions in both regimes. Let $L_z,U_z$
be the corresponding trimmed means. Then a valid sharp interval at that mass is
\[
 [L_1(\pi_{AS}/p_1)-U_0(\pi_{AS}/p_0),
   U_1(\pi_{AS}/p_1)-L_0(\pi_{AS}/p_0)].                \tag{3}
\]
The construction selects each potential survivor subset independently and
couples the selected outcomes into the always-survivor stratum; the other
survival strata fill the remaining masses. Taking the union over positive
feasible masses describes the identified set, or its closure if masses tend
to zero. A zero lower bound allows models with no always survivors, so a
population-wide well-defined target needs a separate positive-mass assumption.
Monotonicity is therefore useful but cannot be justified by an average survival
comparison alone.

\section*{Migration is not death, and three outcomes are joint}
A migrant survivor remains part of the fixed cohort. If outcome observation
indicator $R$ is missing for some survivors, with response rate $r_z$ and
observed mean $m_z$, then
\[
 \E[Y(z)\mid S(z)=1]\in[r_z m_z,r_z m_z+1-r_z].          \tag{4}
\]
Proof: decompose the survivor mean by $R$, and bound the missing conditional
mean in $[0,1]$. This is sharp absent restrictions on missingness. For trimmed
bounds one needs the full survivor distribution, not only (4); optimize over
all completions of its missing mass before trimming. Assuming destination
follow-up is missing at random can narrow this class but must be stated and
supported. Treating migration as a zero outcome changes the estimand.

Schooling, mental health and earnings can each have bounds like (2), yet their
coordinate intervals need not be jointly sharp. The same always-survivor
membership must explain all three observed outcomes. An exact finite-support
implementation assigns weights to observed outcome vectors, constrains selected
mass to $q$, and minimizes any declared linear welfare index. Selecting the
bottom schooling quantiles and the bottom earnings quantiles independently may
require different children and hence an impossible joint construction.

\section*{Remaining causal and intergenerational issues}
Randomization was used to expose the selection problem cleanly. With naturally
occurring exposure, conditional exchangeability or a defensible instrument is
still needed; controlling for post-exposure displacement may remove part of
the total effect. Persistence after parental exposure ends must be distinguished
from repeated childhood exposure. Mediation through parental income, violence
or schooling requires stronger cross-regime assumptions than the total-effect
bounds.

For descendants conceived after the policy, the two regimes can produce different
children. A fixed-cohort individual contrast is then not defined without a
principal-stratum conception rule or a population aggregate target. These bounds
solve a restricted selection problem for existing children and do not bypass
fertility, confounding or long-run measurement. No real-world effect estimate
is asserted.
\begin{thebibliography}{9}
\bibitem{source} S. Tobon, M. M. Sviatschi and N. Cabra-Ruiz (2025),
\emph{Organised Crime}, VoxDevLit 17(1).
\url{https://voxdev.org/voxdevlit/organised-crime}.
The publisher page supplies development and intergenerational context.
The trimming argument is an explicit application of standard selection-bound
reasoning, proved here; it is not a claim of a novel general bounding method.
\end{thebibliography}

\end{document}
